\documentclass[letterpaper, 10 pt, conference]{ieeeconf}

\IEEEoverridecommandlockouts
\overrideIEEEmargins

\usepackage{amsmath,amssymb,amsfonts}
\usepackage{graphicx}
\usepackage{algorithm}
\usepackage{algorithmic}
\usepackage{cite}
\usepackage{xcolor}
\usepackage{siunitx}

\title{RRT* Motion Planning with Control Barrier Function Safety Derived from Poisson Fields}

\author{Maciej Różewicz
    \thanks{Department of Automatic Control and Robotics,
    AGH University of Science and Technology, Kraków, Poland.
    Email: mrozewicz@agh.edu.pl}
}

\begin{document}

\maketitle
\thispagestyle{empty}
\pagestyle{empty}

%%%%%%%%%%%%%%%%%%%%%%%%%%%%%%%%%%%%%%%%%%%%%%%%%%%%%%%%%%%%%%%%%%%%%%%%%%%%%%%%%%%%%%%%%
\begin{abstract}
This paper proposes a~novel framework for safe motion planning that integrates
Control Barrier Function (CBF) safety constraints with the Rapidly-Exploring
Random Tree Star (RRT*) algorithm. The CBF is derived directly from the solution
of a~Poisson equation defined over the free workspace, referred to as
a~\emph{Poisson Safety Function (PSF)}. The PSF encodes obstacle boundaries
and safety gradients in a~smooth, differentiable field that naturally satisfies
Laplace-type spatial regularity. The proposed method, termed
\emph{Safe-CBF-RRT*}, uses the PSF and its spatial derivatives to enforce
CBF-based safety conditions along sampled edges during tree expansion.
Simulation results demonstrate that the method enhances safety and convergence
properties compared to standard RRT*, while maintaining computational tractability.
\end{abstract}

%%%%%%%%%%%%%%%%%%%%%%%%%%%%%%%%%%%%%%%%%%%%%%%%%%%%%%%%%%%%%%%%%%%%%%%%%%%%%%%%%%%%%%%%%
\section{Introduction}
In modern robotics applications, ensuring safe navigation in complex
environments is one of the major challenges. 

Sampling-based motion planners such as RRT and RRT* are widely used for complex
robotic systems due to their scalability and ability to handle nonconvex
constraints. However, they rely primarily on geometric collision checking,
offering no formal guarantees of safety. Control Barrier Functions (CBFs)
provide a~theoretical framework for ensuring safety in continuous-time control
systems by enforcing forward invariance of safe sets.

A~critical challenge in using CBFs for planning is the construction of
a~suitable safety function $h(x)$, which must be continuously differentiable
and satisfy certain regularity conditions. Analytical definitions of $h(x)$
exist for simple geometries, but not for arbitrary obstacle configurations.

In this work, we propose a~method to automatically generate a~smooth,
differentiable safety function using the Poisson equation. The resulting
\emph{Poisson Safety Function (PSF)} is used as the foundation of a~CBF
integrated within RRT*. This combination provides both formal safety
filtering and smooth exploration.

\textbf{Contributions:}
\begin{itemize}
    \item Integration of the derived field into a modified RRT* algorithm.
    \item Empirical results showing improved safety margins and convergence.
\end{itemize}

%%%%%%%%%%%%%%%%%%%%%%%%%%%%%%%%%%%%%%%%%%%%%%%%%%%%%%%%%%%%%%%%%%%%%%%%%%%%%%%%%%%%%%%%%
\section{Related Work}
Sampling-based methods such as RRT and RRT* \cite{karaman2011sampling} have
become standard in high-dimensional motion planning. Extensions incorporating
optimality, dynamics, and risk awareness exist but often lack formal safety guarantees.

Potential field and PDE-based navigation \cite{connolly1990harmonic} have long
been studied for path planning. Harmonic potential fields, governed by Laplace's
equation, ensure smoothness and avoid local minima, yet they do not directly
provide a~safety margin measurable in the CBF framework.

Control Barrier Functions \cite{ames2017control,wang2022safe} ensure set
invariance in control systems but require a~known differentiable safety
function. Existing methods often define $h(x)$ as analytic distances to
obstacles, which may be nondifferentiable in complex maps. Our approach
bridges this gap by generating $h(x)$ through a~Poisson PDE that respects
boundary geometry.


\cite{bahati2025dynamicsafetycomplexenvironments}

%%%%%%%%%%%%%%%%%%%%%%%%%%%%%%%%%%%%%%%%%%%%%%%%%%%%%%%%%%%%%%%%%%%%%%%%%%%%%%%%%%%%%%%%%
\section{Preliminaries}
\subsection{Poisson Safety Function Generation}
Let $\Omega \subset \mathbb{R}^2$ denote the workspace with boundary
$\partial\Omega$ and obstacles $\mathcal{O}_i$. We define a~vector
potential field $\mathbf{u} = [u_x,u_y]^T$ solving
\begin{equation}
    \left\lbrace
    \begin{array}{ll}
        \nabla^2 \mathbf{u} = 0 & \text{in } \Omega \\
        \mathbf{u} = c_i (\mathbf{x}-\mathbf{c}_i) & \text{on } \partial \mathcal{O}_i
    \end{array}
    \right\rbrace
    .
\end{equation}
A scalar field $h(x,y)$, termed the \emph{Poisson Safety Function (PSF)},
is then obtained by solving
\begin{equation}
    \left\lbrace
    \begin{array}{ll}
        \nabla^2 h = -\|\nabla \mathbf{u}\| & \text{in } \Omega \\
        h|_{\partial\Omega}=0 & \text{on } \partial \mathcal{O}_i
    \end{array}
    \right\rbrace
    ,
\end{equation}
ensuring $h>0$ in free space. Gradients $\nabla h = [\partial h/\partial
x,\partial h/\partial y]$ are computed using finite differences.

\subsection{Control Barrier Function Constraint}
For a~path segment between points $a$ and $b$, sampled along $p(s)=a+s(b-a)$,
$s\in[0,1]$, the discrete CBF condition is
\begin{equation}
\nabla h(p) \cdot (b-a) \ge -\kappa h(p)\|b-a\|,
\end{equation}
where $\kappa>0$ controls conservatism. This ensures that motion along
the edge does not approach unsafe regions too rapidly.

%%%%%%%%%%%%%%%%%%%%%%%%%%%%%%%%%%%%%%%%%%%%%%%%%%%%%%%%%%%%%%%%%%%%%%%%%%%%%%%%%%%%%%%%%
\section{Safe CBF-RRT* Algorithm}
Safe CBF-RRT* extends RRT* by filtering edges through CBF validation and
weighting path costs by safety:
\begin{align}
J(a,b) = c L + (1-c)\frac{1}{\bar{h}},
\end{align}
where $L=\|b-a\|$, $\bar{h}$ is mean $h$ along the edge, and $c\in[0,1]$
balances length and safety.

The planner iteratively samples, extends, and rewires as in RRT*, rejecting
edges that violate the CBF condition. A pseudocode outline is shown in
Algorithm~\ref{alg:cbfrrt}.

\begin{algorithm}[t]
\caption{Safe-CBF-RRT*}
\label{alg:cbfrrt}
\begin{algorithmic}[1]
\STATE Initialize tree with $x_{start}$
\FOR{$k=1:N_{iter}$}
  \STATE Sample $x_{rand}$
  \STATE $x_{near} \leftarrow$ Nearest$(x_{rand})$
  \STATE $x_{new} \leftarrow$ Steer$(x_{near}, x_{rand})$
  \IF{CollisionFree($x_{near},x_{new}$) and CBFSafe($x_{near},x_{new}$)}
    \STATE AddNode($x_{new}$); RewireNeighborhood($x_{new}$)
  \ENDIF
  \IF{$x_{new}$ near $x_{goal}$}
    \STATE Connect and return path
  \ENDIF
\ENDFOR
\end{algorithmic}
\end{algorithm}

%%%%%%%%%%%%%%%%%%%%%%%%%%%%%%%%%%%%%%%%%%%%%%%%%%%%%%%%%%%%%%%%%%%%%%%%%%%%%%%%%%%%%%%%%
\section{Results}
\subsection{Simulation Setup}
Experiments were conducted in MATLAB using the PDE Toolbox. The workspace was
a~unit square with three circular obstacles. The grid resolution was
$100\times100$. Parameters used were $\text{stepSize}=1$, $\kappa=20$, $c=0.1$,
and $N_{iter}=5000$.

\subsection{Performance}
Two configurations were tested:
\begin{enumerate}
    \item RRT* without CBF safety (baseline)
    \begin{center}
    \includegraphics[width=1.0\linewidth]{fig/rrt_without_cbf}
    \end{center}
    \item Safe-CBF-RRT* (proposed)
    \begin{center}
    \includegraphics[width=1.0\linewidth]{fig/rrt_with_cbf}
    \end{center}
\end{enumerate}
The PSF produced smooth gradients and continuous safety margins. The proposed
planner consistently generated paths that maintained higher safety levels while
remaining within 10\% of the shortest-path length.

\subsection{Discussion}
The CBF constraint effectively rejected paths approaching obstacles too closely.
The Poisson field provided well-conditioned gradients, improving planner
stability. Despite the additional checks, runtime increased only marginally due
to efficient gridded interpolation.

%%%%%%%%%%%%%%%%%%%%%%%%%%%%%%%%%%%%%%%%%%%%%%%%%%%%%%%%%%%%%%%%%%%%%%%%%%%%%%%%%%%%%%%%%
\section{Conclusion and Future Work}
We presented \emph{Safe-CBF-RRT*}, a motion planning framework combining RRT*
with a~Control Barrier Function derived from a~Poisson Safety Field. The
approach unifies PDE-based potential field generation with formal safety
guarantees. Future work will extend this method to kinodynamic models, uncertain
environments, and real robotic experiments.

%%%%%%%%%%%%%%%%%%%%%%%%%%%%%%%%%%%%%%%%%%%%%%%%%%%%%%%%%%%%%%%%%%%%%%%%%%%%%%%
\bibliographystyle{abbrv}
\bibliography{refs}

\end{document}
